%============================================================
% HALAMAN ABSTRAK  (halaman iii)
% Tulis abstrak dalam Bahasa Indonesia, 300-500 kata.
% Isi: latar belakang (~1 kalimat), permasalahan (1-2 kalimat),
%      solusi (~5 kalimat), kinerja (~2 kalimat).
%============================================================
\cleardoublepage
\chapter*{ABSTRAK}
\addcontentsline{toc}{chapter}{ABSTRAK}
\thispagestyle{fancy}

Aplikasi FoodLAB yang digunakan untuk pemesanan makanan daring di Politeknik Elektronika Negeri Surabaya (PENS) menghadapi tantangan signifikan dalam hal \textit{observability} dan pemulihan insiden pada sisi \textit{backend}-nya. \textit{Observability} merupakan kemampuan suatu sistem untuk memperlihatkan kondisi internalnya melalui data telemetri berupa \textit{metrics}, \textit{log}, dan \textit{traces}, sehingga operator dapat memahami perilaku sistem secara menyeluruh. Sistem \textit{monitoring} konvensional yang berbasis \textit{threshold} statis dan respons manual terbukti tidak memadai untuk mendeteksi anomali secara akurat, sehingga menghasilkan \textit{false-positive rate} yang tinggi serta memperpanjang \textit{Mean Time to Detect} (MTTD) dan \textit{Mean Time to Recovery} (MTTR). Untuk mengatasi permasalahan tersebut, penelitian ini merancang dan mengimplementasikan \textit{framework} \textit{Agentic DevOps Monitoring Agents} (ADMA). ADMA adalah sistem \textit{multi-agent} otonom berbasis kecerdasan buatan yang mampu melakukan \textit{observability}, \textit{root cause analysis}, dan \textit{self-healing} pada infrastruktur \textit{backend} FoodLAB. Sistem ADMA terdiri dari dua komponen utama. Pertama, sebuah Go \textit{binary} yang berjalan di \textit{server} target dan bertugas melakukan \textit{monitoring} secara \textit{real-time} mencakup PostgreSQL, Docker \textit{container} \textit{logs}, dan \textit{resource utilization} sistem. Komponen ini juga melakukan diagnosis berbasis \textit{Large Language Models} (LLM) serta eksekusi remediasi otomatis. Kedua, sebuah \textit{backend} \textit{control plane} berbasis C\# ASP.NET Core yang berfungsi sebagai pusat manajemen, persistensi data, dan antarmuka REST API. Kedua komponen berkomunikasi secara asinkron melalui NATS \textit{messaging}. Hasil implementasi ADMA pada infrastruktur FoodLAB diharapkan dapat meningkatkan kemampuan deteksi anomali secara proaktif, mempercepat proses diagnosis dan pemulihan insiden, serta mengurangi ketergantungan pada intervensi manual dalam pengelolaan sistem.

\vspace{1cm}
\noindent\textbf{Kata Kunci:} \textit{Agentic AI}, \textit{Observability}, \textit{Self-Healing}, \textit{DevOps Monitoring}, FoodLAB, LLM, NATS